% Portable pdfLaTeX setup. Uses standard TeX Live fonts; no font files supplied.
\usepackage[T1]{fontenc}
\usepackage[utf8]{inputenc}
\usepackage[english]{babel}
\usepackage[fontsize=12.2pt]{scrextend}
\usepackage{lmodern}
\usepackage[sc]{mathpazo}
\usepackage{helvet}
\renewcommand{\ttdefault}{lmtt}
\usepackage[a4paper,left=22mm,right=22mm,top=22mm,bottom=22mm,headheight=16pt,headsep=7mm,footskip=12mm]{geometry}
\usepackage{microtype,amsmath,amssymb,xspace}
% Typewriter text must stay plain ASCII: no << >> -- ligatures in code set outside listings.
\DisableLigatures[<,>,-,`,',\,]{encoding = T1, family = tt*}
\usepackage{textcomp}
\usepackage{setspace}
\setstretch{1.055}
\usepackage{xcolor}
% Colour edition (for reading on screen). The grayscale print palette is in the
% Best Practices reader; only the definitions in this block differ.
\definecolor{ink}{HTML}{1B2430}      % text and main headings
\definecolor{accent}{HTML}{12497A}   % deep blue
\definecolor{gold}{HTML}{C8922A}
\definecolor{muted}{HTML}{5A6675}
\definecolor{pale}{HTML}{EEF3F8}     % light blue tint
\definecolor{warm}{HTML}{FBF4E6}
\definecolor{mint}{HTML}{FFFFFF}
\definecolor{okgreen}{HTML}{2E7D5B}
\definecolor{palegreen}{HTML}{EDF6F1}
\definecolor{alert}{HTML}{B5483A}
\definecolor{palealert}{HTML}{FBEFEC}
\definecolor{paleslate}{HTML}{F1F3F5}
\definecolor{codebg}{HTML}{F7F8FA}
\definecolor{codeframe}{HTML}{D5DAE0}
\definecolor{codegreen}{HTML}{4F7A5A}
\definecolor{coderust}{HTML}{A8452E}
\definecolor{navy}{HTML}{0E2A47}
\usepackage{graphicx,float,rotating}
\usepackage{xurl,ragged2e}
\definecolor{mint}{gray}{1}
\usepackage{tikz,pgfplots}
\usetikzlibrary{arrows.meta,positioning,calc,shapes.geometric,fit,backgrounds}
\pgfplotsset{compat=1.18}
\usepackage{booktabs,longtable,tabularx,array}
\renewcommand{\tabularxcolumn}[1]{>{\RaggedRight\arraybackslash}p{#1}}
\usepackage[font=small,labelfont={sf,bf,color=ink},labelsep=period,skip=8pt]{caption}
\usepackage{enumitem}
\setlist{leftmargin=*,topsep=6pt,itemsep=5pt,parsep=2pt}
\usepackage{titlesec}
\newcommand{\BookChapterLabel}{CHAPTER}
\titleformat{\chapter}[display]
 {\sffamily\bfseries\color{ink}\raggedright\hyphenpenalty=10000\exhyphenpenalty=10000}
 {\color{gold!85!black}\fontsize{13}{17}\selectfont \BookChapterLabel\ \thechapter}
 {10pt}{\fontsize{25}{29}\selectfont}[\vspace{7pt}{\color{gold}\rule{30mm}{2.2pt}}]
\titlespacing*{\chapter}{0pt}{4pt}{22pt}
\titleformat{\section}{\sffamily\bfseries\Large\color{accent}}{\thesection}{0.65em}{}
\titlespacing*{\section}{0pt}{17pt plus 3pt}{7pt}
\titleformat{\subsection}{\sffamily\bfseries\large\color{ink}}{\thesubsection}{0.65em}{}
\titlespacing*{\subsection}{0pt}{12pt plus 2pt}{5pt}
\usepackage{fancyhdr}
\pagestyle{fancy}
\fancyhf{}
\fancyhead[L]{\footnotesize\sffamily\color{muted}\nouppercase{\leftmark}}
\fancyfoot[L]{\footnotesize\sffamily\color{muted}\BookShortTitle}
\fancyfoot[R]{\sffamily\bfseries\color{accent}\thepage}
\renewcommand{\headrulewidth}{0.6pt}\renewcommand{\headrule}{{\color{gold}\hrule width\headwidth height\headrulewidth}}
\renewcommand{\chaptermark}[1]{\markboth{\thechapter\quad #1}{}}
\fancypagestyle{plain}{\fancyhf{}\fancyfoot[L]{\footnotesize\sffamily\color{muted}\BookShortTitle}\fancyfoot[R]{\sffamily\bfseries\color{accent}\thepage}\renewcommand{\headrulewidth}{0pt}}
\usepackage[most]{tcolorbox}
\tcbset{enhanced,breakable,arc=0.8mm,boxrule=0.35pt,colframe=accent!35,left=10pt,right=10pt,top=8pt,bottom=8pt,before skip=10pt,after skip=11pt}
\newtcolorbox{researchquestion}{colback=pale,borderline west={1.5pt}{0pt}{accent},fontupper=\itshape}
\newtcolorbox{sourcepanel}{colback=black!3,fontupper=\small}
\newtcolorbox{practicebox}[1]{colback=mint,colframe=accent!35,title={\IfBeginWith{#1}{Practice}{#1}{Practice: #1}},colbacktitle=mint,coltitle=ink,fonttitle=\sffamily\bfseries,borderline west={1.5pt}{0pt}{accent},titlerule=0pt}
\newtcolorbox{readerroute}{colback=pale,borderline west={1.5pt}{0pt}{accent},fontupper=\small,title={Reading route},colbacktitle=pale,coltitle=ink,fonttitle=\sffamily\bfseries,titlerule=0pt}
\newtcolorbox{workedexample}[1]{colback=mint,colframe=accent!35,title={\IfBeginWith{#1}{Worked example}{#1}{Worked example: #1}},colbacktitle=mint,coltitle=ink,fonttitle=\sffamily\bfseries,borderline west={1.5pt}{0pt}{accent},titlerule=0pt}
\usepackage{listings}
\lstset{basicstyle=\small\ttfamily,breaklines=true,columns=fullflexible,keepspaces=true,showstringspaces=false,frame=single,rulecolor=\color{black!20},xleftmargin=5pt,xrightmargin=5pt,aboveskip=10pt,belowskip=10pt}
\usepackage{needspace}
\usepackage{etoolbox}
\pretocmd{\section}{\Needspace{8\baselineskip}}{}{\PackageError{booksetup}{Cannot protect section headings}{}}
\pretocmd{\subsection}{\Needspace{5\baselineskip}}{}{\PackageError{booksetup}{Cannot protect subsection headings}{}}
\usepackage{csquotes,xstring}
\usepackage[authoryear,round]{natbib}
\let\parencite\citep
\let\textcite\citet
\usepackage{makeidx}
\makeindex

\usepackage[colorlinks=true,linkcolor=accent,citecolor=accent,urlcolor=accent,unicode=true,bookmarksnumbered=true]{hyperref}
\usepackage{bookmark}
\hypersetup{pdftitle={\BookTitle: \BookSubtitle},pdfauthor={\BookAuthor},pdfsubject={ENG 101 Programming for Engineers, reading notes},pdfkeywords={programming languages, engineering, Python, course notes}}
\setcounter{tocdepth}{0}
\setcounter{secnumdepth}{1}
\setlength{\parindent}{0pt}
\setlength{\parskip}{5pt plus 1pt}
\setlength{\emergencystretch}{3em}
\renewcommand{\arraystretch}{1.20}
\setlength{\tabcolsep}{5pt}
\setlength{\textfloatsep}{15pt plus 3pt minus 2pt}
\setlength{\intextsep}{13pt plus 3pt minus 2pt}
\clubpenalty=10000
\widowpenalty=10000
\raggedbottom
% Larger secondary text than the default 12pt book class, especially figure labels.
\makeatletter
\renewcommand\small{\@setfontsize\small{10.8}{13.2}}
\renewcommand\footnotesize{\@setfontsize\footnotesize{10}{12.2}}
\renewcommand\scriptsize{\@setfontsize\scriptsize{9}{10.7}}
\makeatother
\newcommand{\Field}[2]{\par\noindent\textbf{#1}\ #2\par}
\newcommand{\OnePageHeading}[1]{\section*{#1}}
\newcommand{\ProjectChapter}[2]{\chapter*{#1}\markboth{#1}{}\phantomsection\addcontentsline{toc}{chapter}{#1}\label{#2}}
\newcommand{\PlanSection}[2]{\section{#1}#2\par}
\newcommand{\BookPart}[3]{%
 \clearpage\refstepcounter{part}\markboth{#1}{}%
 \addcontentsline{toc}{part}{\thepart\space #1}\thispagestyle{plain}%
 \vspace*{22mm}{\sffamily\bfseries\color{accent}\Large PART \thepart\par}%
 \vspace{9mm}{\sffamily\bfseries\color{ink}\fontsize{31}{37}\selectfont\raggedright\hyphenpenalty=10000\exhyphenpenalty=10000 #1\par}%
 \vspace{9mm}{\color{accent}\rule{\textwidth}{1pt}\par}%
 \vspace{6mm}{\large #2\par}\vspace{10mm}#3\par\vfill\clearpage}
\newcounter{casebrief}
\renewcommand{\thecasebrief}{\Alph{casebrief}}
\newcommand{\CaseChapter}[1]{\clearpage\refstepcounter{casebrief}\chapter*{Case \thecasebrief: #1}\markboth{Case \thecasebrief: #1}{}\addcontentsline{toc}{chapter}{Case \thecasebrief: #1}\label{case:\thecasebrief}}
\newcommand{\CaseRef}[1]{Case~\ref{case:#1}}
\newcommand{\CheckField}[2]{\textbf{#1} & #2\\\addlinespace[5pt]}
\newcommand{\WritingPrompt}[1]{\begin{sourcepanel}\textbf{Writing prompt.} #1\end{sourcepanel}}

% Full titles in Contents/bookmarks, compact titles only in running heads.
\newcommand{\BookChapter}[2][]{\chapter{#2}\if\relax\detokenize{#1}\relax\else\chaptermark{#1}\fi}
\newcommand{\ExerciseLevel}[1]{{\sffamily\small\color{accent}[#1]}\enspace}
